\subsection{The next two ordinary obstruction degrees}
\label{sec:kt-quartic-quintic-obstruction}

Retain (CM.1)--(CM.4), and now assume all three auxiliary units
at exponents \(0,e,2e\).  The hypotheses imply \(p\ge7\):
\(p=5\) and even \(e\) would give \((p-1)\mid2e\).
Thus exponential and logarithmic coordinates through degree five
are defined.  Write \(A_j=k[s,t]/(s,t)^{j+1}\),
\(F=\mathbf Q(\mu_p)\), \(K=EF\), and
\(A_J=\operatorname{Cl}(J)\{p\}\).
All assertions concern the specified equivariant ordinary Artin
families.  They do not assert that a characteristic-zero source
point supplies such a family.

Set
\[
 P=E_{13},\qquad H=\operatorname{diag}(1,0,-1),\qquad Q=E_{31}.
                                                        \tag{HQ.1}
\]
The permitted homogeneous terms after (CM.5) are exactly
\[
 \begin{array}{c|c|c|c}
 \text{degree}&\text{monomial}&\text{Lie vector}&\text{character}\\\hline
 4&s^3t&P&\eta\psi^2\\
 4&s^2t^2&H&1\\
 4&st^3&Q&\eta\psi^{-2}\\
 5&s^3t^2&X&\eta\psi\\
 5&s^2t^3&Y&\psi^{-1}
 \end{array}                                             \tag{HQ.2}
\]
Indeed the integer torus weights of \(\mathfrak{sl}_3\)
are \(-2,-1,0,1,2\).  The involution signs of
\(P,H,Q,X,Y\) are respectively \(-,+,-,+,-\).
These weights and signs give (HQ.2); there are no additional
permitted monomials.  In particular the positive long root
at degree four cannot be omitted.

\begin{theorem}[KT-ORDINARY-QUARTIC-CONDITION]
Fix an ordinary equivariant lift of the quadratic family to
\(A_3\).  It has a well-defined remaining quartic obstruction
\[
 \mathfrak o_4^+\in H^2(G_{\mathbf Q,S},k(\eta\psi^2)).
                                                        \tag{HQ.3}
\]
The cubic family lifts to \(A_4\), with the same ordinary
and away-prime conditions, if and only if
\(\mathfrak o_4^+=0\).  The other two quartic directions
have zero relative obstruction spaces.

The target of (HQ.3) has the exact description
\[
 H^2(G_{\mathbf Q,S},k(\eta\psi^2))
 \simeq
 \operatorname{Hom}_{\operatorname{Gal}(K/\mathbf Q)}
       (A_K/pA_K,k(\eta\omega^{1-2e}))^\vee.
                                                        \tag{HQ.4}
\]
Here the dual is over \(k\), using the trace pairing.
The extra unit
\[
                       B_{1,\eta\omega^{-2e}}\in\mathbf Z_p^\times
                                                        \tag{HQ.5}
\]
is sufficient to make this target zero.  It is not one of
the three assumed units in general.  This does not claim
that (HQ.5) is necessary for the specific class (HQ.3)
to vanish.
\end{theorem}

\begin{proof}
Use \(A,B,C,a,b,c\) as in (CM.9), and additionally put
\[
 M=m(g),\quad N=n(g),\quad
 m=\psi(g)m(h),\quad n=\eta(g)\psi(g)^{-1}n(h),
 \qquad f=Ab-Ba.
                                                        \tag{HQ.6}
\]
The repeated letters \(m,n\) on the right denote the
transported values at \(h\), only within this calculation.
The degree-four BCH defects, in the order \(P,H,Q\), are
\[
\begin{split}
 \Omega_{4,+}
   & =aM-Am+\tfrac12(A-a)(Ac-aC)-\tfrac14Aa f,\\
 \Omega_{4,0}
   & =\tfrac12(An-Bm+bM-aN)+\tfrac12(ABc+abC)\\
   &\hspace{12mm}-\tfrac14(C+c)(Ab+Ba)-\tfrac18(Ab+Ba)f,\\
 \Omega_{4,-}
   & =bN-Bn-\tfrac12(B-b)(Bc-bC)+\tfrac14Bb f.
\end{split}                                               \tag{HQ.7}
\]
If the new one-cochains are \(r_+,r_0,r_-\), their
equations in the convention of (CM.14) are
\(dr_i=-\Omega_{4,i}\).
One obtains (HQ.7) from the terms of total degree four in
\[
 A+B+\tfrac12[A,B]+\tfrac1{12}[A,[A,B]]
 +\tfrac1{12}[B,[B,A]]-\tfrac1{24}[B,[A,[A,B]]],
                                                        \tag{HQ.8}
\]
where in (HQ.8) the letters \(A,B\) denote the two
entire logarithmic matrices.  Alternatively the finite
matrix exponential, product, and logarithm give the same
three expressions directly.  The independent program
\texttt{quartic\_bch\_check.py} uses this latter computation
and checks all five degree-four monomials.
Associativity makes these defects two-cocycles: the
degree-four kernel is an additive central layer over
the fixed residual action.

We justify the actual local spaces before using global
cohomology.  At \(p\), the checked formally smooth
ordinary functor in KT-ORDINARY-INTEGRAL-COMPONENT gives
local lifts.  Its proof applies to these small extensions;
the even residual exponent differences never equal the
cyclotomic exponent one.  Fixed determinant is obtained
by removing the trace, since \(3\) is a unit.
The permitted local differences are all \(H^1\) in
an upper root, the unramified \(H^1\) in the diagonal
direction, and zero \(H^1\) in a lower root.  The last
assertion retains flag conjugation: a lower-root
difference can be produced by changing the ordinary
flag precisely when it is a local coboundary.

At \(q\mid D\), the quotient obtained by killing
inertia over \(E\) has an inertia subgroup of order
at most two and a procyclic unramified quotient.
It has \(p\)-cohomological dimension one.  Hence
acceptable local lifts exist across the square-zero
layer.  Its tangent \(H^1\) is zero for an
\(\eta\)-ramified coefficient and is ordinary
unramified \(H^1\) for a rational coefficient.
Projection to the required integer weight and
involution sign preserves these linear lifting
equations.  This gives equivariant local lifts as
well; it is not an inference from raw global
\(H^2\) alone.

For \(k(\eta\psi^{-2})\), the unit at \(2e\)
annihilates the even mirror
\(A_K^{\eta\omega^{2e+1}}\) and its strict
negative group by KT-AUXILIARY-UNIT-SELMER-BRIDGE.
Exactly the proof of (CC.1), using only
\(2e\not\equiv0\pmod{p-1}\), gives
\[
 H^2(G_{\mathbf Q,S},k(\eta\psi^{-2}))=0,
 \quad H^1(G_{\mathbf Q,S},k(\eta\psi^{-2}))
      \xrightarrow{\sim}H^1(\mathbf Q_p,k(\eta\psi^{-2})).
                                                        \tag{HQ.9}
\]
At every \(q\mid D\) its local \(H^1\) is zero.
Consequently a global solution can be corrected to
match the required lower-root local coboundary at
\(p\); the negative relative obstruction space is zero.

For the trivial diagonal coefficient, unramified
conditions away from \(p\) reduce the relative problem
to \(G_{\mathbf Q,p}\) by the same excision as (CE.3).
Ray class theory over \(\mathbf Q\), or Kronecker--Weber,
gives
\[
 \dim_k H^1(G_{\mathbf Q,p},k)=1,\qquad
 H^1(G_{\mathbf Q,p},k)
    \xrightarrow{\sim}
 H^1(\mathbf Q_p,k)/H^1_{\rm unr}(\mathbf Q_p,k).
                                                        \tag{HQ.10}
\]
The character of the degree-\(p\) subfield of
\(\mathbf Q(\mu_{p^2})\) generates the first group
and is ramified at \(p\).  The target also has
dimension one by local Euler characteristic.
Global Euler characteristic for the trivial even
module gives \(1-1+h^2=0\), so
\(H^2(G_{\mathbf Q,p},k)=0\).
The relative exact sequence and (HQ.10) show that
the diagonal obstruction space is zero.

For the positive coefficient \(M=k(\eta\psi^2)\),
its Cartier dual is \(k(\eta\omega^{1-2e})\).
Every local \(H^2\) in \(S\) is zero: the dual
has nontrivial inertia at \(p\), and quadratic
ramification at \(q\mid D\).  The local \(H^1\)
at such \(q\) is also zero.  Since the positive
local space at \(p\) is all \(H^1\), the relative
obstruction is exactly raw global \(H^2\).
Set \(\mathfrak o_4^+=-[\Omega_{4,+}]\).
The independent central directions and the preceding
two zero relative groups prove the lifting criterion.

Poitou--Tate duality identifies this global \(H^2\)
with the dual of strict \(H^1\) in the displayed
Cartier-dual character.  Restriction to \(K\),
prime-to-\(p\) averaging, and class field theory
identify that strict group with the Hom space in
(HQ.4).  There is no prime-class quotient: at a
prime over \(p\) the cyclotomic inertia exponent
\(1-2e\) is nontrivial, and at a prime over \(D\)
the quadratic inertia character is nontrivial.
The permutation modules of these primes therefore
have no component in that character.  These same
facts show that every class-group homomorphism in
(HQ.4) satisfies the required strict local conditions.
Finally apply the auxiliary-unit bridge with the
nonzero even residue \(a\equiv-2e\pmod{p-1}\).
Its even mirror is \(\eta\omega^{1-2e}\), proving
the sufficient condition (HQ.5).
\end{proof}

\begin{theorem}[KT-QUINTIC-RELATIVE-TARGET]
Suppose the fixed cubic family admits a quartic ordinary lift.
For any such lift the two possible quintic obstructions have
coefficients \(k(\eta\psi)\) and \(k(\psi^{-1})\).
The upper relative obstruction space is
\[
 H^2(G_{\mathbf Q,S},k(\eta\psi))
 \simeq\operatorname{Hom}_{\operatorname{Gal}(K/\mathbf Q)}
       (A_K/pA_K,k(\eta\omega^{1-e}))^\vee.
                                                        \tag{HQ.11}
\]
The additional unit \(B_{1,\eta\omega^{-e}}\) is
sufficient to kill (HQ.11).

Let \(\mathcal O^-_5\) denote the rational lower
relative obstruction space.  Its local condition at
\(p\) is zero, and its away-prime conditions are
unramified.  Put
\[
 d=\dim_k\operatorname{Hom}_{\operatorname{Gal}(F/\mathbf Q)}
                       (A_F/pA_F,k(\omega^{-e})).
                                                        \tag{HQ.12}
\]
There is an exact sequence
\[
\begin{split}
0\longrightarrow&
 \operatorname{coker}\bigl(H^1(G_{\mathbf Q,p},k(\psi^{-1}))
                  \to H^1(\mathbf Q_p,k(\psi^{-1}))\bigr)\\
\longrightarrow&\mathcal O^-_5
\longrightarrow H^2(G_{\mathbf Q,p},k(\psi^{-1}))
\longrightarrow0,
\end{split}                                               \tag{HQ.13}
\]
and precisely
\[
                         \dim_k\mathcal O^-_5=d+1.
                                                        \tag{HQ.14}
\]
The dimension \(d\) is that of the original residual
strict rational negative Selmer group.  The raw last
term in (HQ.13) has the odd mirror character
\(\omega^{e+1}\); it cannot replace the full relative
space.

For the fixed cubic family, its rational lower
quintic obstruction is independent of the choice
of acceptable quartic lift.  If the upper space
(HQ.11) is zero, the quartic family lifts through
degree five if and only if that rational relative
class is zero.  No nonzero value of this class is
asserted.
\end{theorem}

\begin{proof}
Table (HQ.2) lists the entire kernel of
\(A_5\to A_4\) allowed by the grading and involution.
It is an additive central layer, so the homomorphism
defect of any set-theoretic lift gives the two
two-cocycles, with the same relative construction
as (CM.14).  The local lifting and flag argument
in the preceding proof applies without change.
For the upper character the local space at \(p\)
is full and the spaces at \(q\mid D\) vanish.
The argument for (HQ.4), with \(2e\) replaced by
\(e\), proves (HQ.11) and its sufficient opposite
unit condition.

For the lower rational character, a permitted
local difference at \(p\) is a coboundary: it
comes from changing the ordered ordinary flag.
Thus \(L_p=0\), not full local \(H^1\).
At \(q\mid D\), killing inertia over \(E\)
gives exactly unramified \(H^1\) by the
order-at-most-two averaging already proved.
Apply etale excision at those primes while
retaining the \(p\)-primitive.  The result is
the mapping-fibre degree-two group for
\(R\Gamma(G_{\mathbf Q,p},k(\psi^{-1}))\to
R\Gamma(\mathbf Q_p,k(\psi^{-1}))\).
Local \(H^2\) is zero, because its dual inertia
character \(\omega^{e+1}\) is nontrivial.
The mapping-fibre sequence is exactly (HQ.13).

Write \(r=\dim_k H^2(G_{\mathbf Q,p},k(\psi^{-1}))\).
The Kummer and class-group argument in (CE.4)--(CE.5)
identifies this group with the scalar extension of
\((A_F/pA_F)^{\omega^{e+1}}\).
Since \(\psi^{-1}\) is even and nontrivial,
global Euler characteristic gives
\(\dim H^1(G_{\mathbf Q,p},k(\psi^{-1}))=r\).
Local \(H^1\) has dimension one: both its
local \(H^0\) and \(H^2\) vanish.

The kernel of that localization has dimension
\(d\).  To see the exact identification, restrict
to \(F\).  A strict class becomes an everywhere
unramified homomorphism.  Conversely an equivariant
homomorphism from \(A_F/pA_F\) in character
\(\omega^{-e}\) gives a global class unramified
outside \(p\); its restriction at the prime above
\(p\) is zero because that prime is principal.
The local and global restriction maps are
injective by averaging the degree \(p-1\).
This proves the claimed equality with the
original residual strict group.  Hence the
cokernel in (HQ.13) has dimension \(1-(r-d)\),
and adding \(r\) proves (HQ.14).

Finally compare two ordinary quartic lifts of
the same cubic family.  Their difference is
a one-cocycle in the three lines of degree four.
In the negative line its global strict class
is zero by the unit at \(2e\); in the diagonal
line its everywhere-unramified class is zero
by the class group of \(\mathbf Q\).
Modulo permitted global conjugations, only
the positive \(s^3tP\) difference remains.
Its brackets with the degree-one terms can
affect only \(s^3t^2X\), since
\([P,X]=0\) and \([Y,P]=X\).
They cannot affect \(s^2t^3Y\).
Thus the rational lower relative obstruction
does not change.  The two central degree-five
directions are independent at this order;
if (HQ.11) vanishes, solving the remaining
relative equation is therefore necessary
and sufficient for the degree-five lift.
\end{proof}

\paragraph{The exact remaining requirements.}
The three given auxiliary units remove the negative
quartic direction and the diagonal direction, but
do not in general remove the positive quartic
target (HQ.4).  An opposite auxiliary unit suffices
there; a second opposite unit suffices for the
upper quintic target.  These are sufficient group
vanishing inputs, not necessary conditions on
individual obstruction values.  The rational
quintic condition then sees a relative space of
dimension one plus the original target dimension.
No injectivity, nonzero value, or general lifting
contradiction follows from this dimension count.
The actual source-to-jet conditions remain
separate.  This section supplies exact later
lifting targets and a quartic cochain formula;
it does not settle an additional arithmetic gate.

\paragraph{Inputs and verification scope.}
The duality, Euler, ordinary-component, and auxiliary-unit
inputs are the same checked statements used in
(CM.12)--(CM.15), (CC.1)--(CC.3), and (AB.3)--(AB.5).
The diagonal calculation additionally uses the ray class
description over \(\mathbf Q\): Milne,
\emph{Class Field Theory}, Theorem~0.5 and the full
conductor paragraph, printed pp.~5--6, including
\(C_{m\infty}\simeq(\mathbf Z/m\mathbf Z)^\times\).
That exact statement and its surrounding conductor
scope were read; global class field theory is imported.
The independent rational matrix-exponential program
checks (HQ.7) and the six displayed bracket identities
used to identify the permitted directions.
No arithmetic norm or Massey value is computed.
